\documentclass{article}
\usepackage{propositions}

%% Which mode an environment sets its items in, across nesting.
%%
%% inlineprop opens no list, so its items are set in running text -- and that
%% has to hold wherever the environment sits.  The mode used to be inferred
%% from the nesting level on the way out of a prop (display if the level was
%% still above zero), which reads every enclosing environment as a prop; an
%% inlineprop is not one.  So an inlineprop inside a prop set its items as
%% list items of that prop, and an inlineprop went back to setting list items
%% after a prop nested inside it closed.  Each environment now remembers the
%% mode it interrupted and gives that back.
%%
%% Also: a bare \item inside an inlineprop supplies its own interword space,
%% the scan for its optional argument having eaten the source's.  \item[] is
%% the way to ask for none.

\begin{document}

\section{The mode across nesting}

\noindent\rule{\textwidth}{0.3pt}
\begin{prop}
  \item An ordinary displayed item. \label{m:1}
  Now an inlineprop: \begin{inlineprop}\item alpha. \item beta.\end{inlineprop}
  Those two must run on inside this paragraph, lettered, not broken out onto
  lines of their own.
  \item A second displayed item. \label{m:2}
\end{prop}
\noindent\rule{\textwidth}{0.3pt}

\begin{inlineprop}
  \item first, inline. \label{m:3}
  \begin{prop}\item A displayed item nested inside the inlineprop.\end{prop}
  \item second, which must still be inline. \label{m:4}
\end{inlineprop}

\noindent\rule{\textwidth}{0.3pt}

An inlineprop within an inlineprop stays inline throughout:
\begin{inlineprop}\item outer. \begin{inlineprop}\item inner one.
\item inner two.\end{inlineprop} back in the outer one.\end{inlineprop}

A prop within a prop is unaffected:
\begin{prop}
  \item One.
  \begin{prop}\item A sub-item.\end{prop}
  \item Two.
\end{prop}
\noindent\rule{\textwidth}{0.3pt}

A footnote interrupts all of it.\footnote{Inline here:
\begin{inlineprop}\item x.\end{inlineprop} and displayed here:
\begin{prop}\item y.\end{prop}}
\begin{prop}
  \item The item after the footnote must still be a list item.
\end{prop}
\noindent\rule{\textwidth}{0.3pt}

References: [\ref{m:1}] [\ref{m:2}] [\ref{m:3}] [\ref{m:4}].

\section{The space after a bare \texttt{\string\item}}

\begin{inlineprop}
  \item Bare: the label must be followed by a space. \item[P] With an
  argument: the source's space, as before. \item[]Empty argument: no space,
  which is how to ask for none.
\end{inlineprop}

\end{document}
